\section{The complete deterministic-specialization calculus}
\label{sec:deterministic-specialization}

\subsection{Syntax, computation, and all primitive rules}

There is one numerical sort.  Start from finite raw terms
\[
 t ::= x\mid 0\mid S(t)\mid I(n,a;u.h).
\]
Only u in h is bound by the iterator; n and a are outside that binder.
The zero-normalization map acts recursively on subterms and contracts
\(I(0,a;u.h)\) to a, but does not contract successor-counter iterators.
The terms of the present calculus are the fixed points of this map.
Every substitution is capture avoiding and followed by this normalization,
including a substitution into an iterator counter that newly becomes zero.
For readability, expressions below always denote their zero-normal forms.

Formulas and clauses have grammars
\[
 F ::= (s=t)\mid D_{x,y}F,\qquad [A;y.B(y)]\quad(y\notin\mathrm{FV}(A)).
\]
Both x and y are bound by D.  A context is a finite ordered list of clause
occurrences; occurrences remain distinct even when their contents agree.
An ordinary assumption A is represented by the constant clause
\([0=0;y.A]\), with y unused.  Put \(1=S0\), and designate \(0=1\)
as contradiction.  Clause order is used only to choose a specialization
target; clauses are reusable, and no assumption is consumed by a rule.

The following five schemas are the entire primitive inventory.  There
are no additional logical, equality, computation, successor, induction,
weakening, exchange, or cut rules.
\[
 \frac{}{\Gamma\vdash t=t}\ (\mathsf{Refl}),\qquad
 \frac{\Gamma\vdash I(c,a;u.h)=t\qquad
       \Gamma\vdash I(Sc,a;u.h)=r}
      {\Gamma\vdash r=h[t/u]}\ (\mathsf{Step}).
\]
In Step the first iterator is normalized, so it is just a when c is zero.
Repeated expressions must be syntactically identical; no hidden equality
decision or evaluation oracle occurs in this rule.
\[
 \frac{\Gamma\vdash D_{x,y}B(x,y)\qquad
       [B(b,0);y.B(b,y)],\Gamma\vdash F}
      {\Gamma\vdash F}\ (\mathsf{DE}).
\]
Here b is fresh for \(\Gamma,F\) and all passive matrix parameters.
\[
 \frac{[0=0;v.B(t,\eta)],\Gamma\vdash B(t,Sz)}
      {\Gamma\vdash D_{x,y}B(x,y)}\ (\mathsf{DI}),\qquad
 \eta=\begin{cases}
 z,&B\text{ is a literal equation and }t\text{ is literally }0,\\
 0,&\text{otherwise}.
 \end{cases}
\]
Here v is unused, and z is fresh for \(\Gamma,t\) and all passive matrix
parameters.  An equation beneath a D is not a literal equation.  The
fifth schema is Use$_*$ below.  The witness-selected DI convention and
its recovery of the missing equation-case instances at witness zero are
inherited from the existing uniform-Step calculus; they are not altered
by the present specialization reduction.

Store context
clauses newest first.  For a nonempty parameter scope let \(b_*\) be its
last parameter; for the empty scope put \(b_*=0\).  Replace unrestricted
clause use by the following rule:
\[
 \frac{\Gamma\vdash A}{\Gamma\vdash B(v_i)}
 \quad [A;y.B(y)]=\Gamma_i,
 \qquad v_i=\begin{cases}1&i=0,\\ b_*&i>0.\end{cases}
 \tag{Use$_*$}
\]
Thus the newest clause is used at \(1=S0\); every older clause is used
at the current last parameter.  The certificate supplies its clause
index and antecedent proof, with no numerical target or mode bit.

\subsection{The arithmetic transport used in the proof}

For \(J_s=I(1,s;u.u)\), Step sends premises \(s=t\) and \(J_s=r\)
to \(r=t\).  If \(p:s=t\) and \(q:s=r\), then
\[
 \mathsf{Step}\bigl(p,\mathsf{Step}(q,\mathsf{Refl}(J_s))\bigr)
       :r=t.
 \tag{Pivot}
\]
This gives symmetry by putting the first premise equal to \(s=s\), and
then transitivity.  For a term template h, put \(K=I(1,s;u.h)\).
Step on reflexivities gives \(K=h(s)\); Step on \(p:s=t\) and that
equation gives \(h(s)=h(t)\).  Hence transport in an equation template
\(F(v)=(a(v)=b(v))\) follows from
\[
 a(t)=a(s),\qquad a(s)=b(s),\qquad b(s)=b(t).
\]
This derivation uses only Refl and Step and opens no new scope.
Applying Step to the two appropriate reflexivities also gives successor
computation.  In particular \(C(0,a,b)=a\) is definitional and
\(C(Sz,a,b)=b\) is derivable for the selector
\(C(y,a,b)=I(y,a;u.b)\).  These observations prove every arithmetic
operation needed by the specialization construction independently of
the logical rules.

\subsection{Exact consequence equivalence}

\begin{theorem}[Deterministic-specialization equivalence]
The resulting calculus derives exactly the same open sequents as the
unrestricted five-rule clause calculus.  The ordinary numerical and
formula interpretations are unchanged.
\end{theorem}

\begin{proof}
Soundness is immediate because every new use is an old use.  We give a
constructive completeness proof.  Scope extension is not assumed: it
changes the last parameter, and inserting a clause changes which clause
is newest.  Those two effects are handled by the proof-producing
procedures below.

Let \(C(y,a,b)=I(y,a;u.b)\), and let
\[
 N_t=D_{x,y}\bigl(C(y,0,x)=C(y,0,t)\bigr).
 \tag{Name}
\]
Introduce \(N_t\) at witness \(t\).  Its required successor equation
has literally equal sides, so reflexivity proves the branch in either
introduction mode.  Eliminate \(N_t\) and call the fresh parameter
\(b\).  The newly inserted clause is
\[
 \mathcal N_{b,t}=[0=0;\ y.C(y,0,b)=C(y,0,t)].
 \tag{N}
\]
Since it is newest, Use$_*$ applies it at 1.  Reflexivity supplies its
antecedent, and selector computation gives \(b=t\).  This particular
proof is valid immediately after insertion; we do not pretend that it
automatically remains valid after further context extensions.

A \emph{persistent premise procedure} is a procedure that produces a
proof of its conclusion after every injective renaming of its free
parameters and every embedding of its retained clause occurrences into
a larger context and scope.  This definition describes an external
finite proof-producing algorithm, not an added structural inference
rule.  A retained constant clause for an ordinary assumption is
persistent: either permitted target gives the same formula.  Refl and
arithmetic combinations of persistent premises are also persistent.

\paragraph{Atomic specialization comes first.}
Let \([A;y.E(y)]\) be any retained clause whose body E is a literal
equation, let P be a persistent procedure for A, and let t be any term.
To specialize E at t, introduce and eliminate \(N_t\).  In its branch,
request P in the enlarged context, and apply the original clause at the
new last parameter b.  The original clause is now older, so this is
exactly its permitted use.  The result is \(E(b)\).  The immediately
available proof \(b=t\) from (N), together with equational congruence,
gives \(E(t)\).  Close the elimination.  Call this procedure \(U_0\).
The arithmetic transport here opens no scope, so it does not require a
persistent proof of \(b=t\).  The whole recipe is uniform under context
and scope embeddings, so \(U_0\) produces a persistent output whenever
P is persistent.

In particular, after (N) has ceased to be newest, its equation \(b=t\)
still has a persistent proof: use \(U_0\) to specialize that retained
atomic clause at 1, with Refl of zero as its persistent antecedent, and
apply selector computation.  This use of \(U_0\) is not circular.
Its own auxiliary naming clause is newest when used at 1, and its final
transport is atomic.  This is the essential extra step that permits the
deterministic protocol.

\paragraph{Mutual recursion by formula rank.}
Rank is the number of D binders above the final equation.  We construct
\(T_r\) and \(U_r\) for every finite rank r:
\begin{itemize}
\item \(T_r\) transports a formula F of rank r along a persistent proof
of \(s=t\), from an arbitrary supplied proof of \(F(s)\) to a proof of
\(F(t)\).  The supplied proof of \(F(s)\) need not be persistent.
\item \(U_r\) specializes a retained clause \([A;y.F(y)]\), with F of
rank r, at an arbitrary term t, provided its antecedent A is supplied by
a persistent procedure.  It produces a persistent proof procedure.
\end{itemize}
For \(T_0\), use the equational congruence already derived from Step;
\(U_0\) was constructed above.

Assume the operations are known below rank r, and write
\(F(v)=D_{x,y}G(v,x,y)\).  For \(T_r\), use the arbitrary supplied
proof of \(F(s)\) immediately as the major premise of DE.  Let its fresh
witness be c.  The context now contains
\([G(s,c,0);y.G(s,c,y)]\).  Introduce \(F(t)\) at witness c.
Because c is a fresh parameter, it is not literally zero; this is a case
introduction, including when G is an equation.  The local assumption
\(G(t,c,0)\) is a constant clause and therefore persistent.  Transport
it back along the persistent equality by \(T_{r-1}\), obtaining a
persistent proof of \(G(s,c,0)\).  Use \(U_{r-1}\) to specialize the
retained closure at \(Sz\), then transport forward by \(T_{r-1}\) to
obtain \(G(t,c,Sz)\).  Close DI and DE.  The original arbitrary minor
proof was used outside the new scope and was never weakened.  The recipe
is persistent when that minor premise is supplied persistently.

Now construct \(U_r\).  Introduce and eliminate \(N_t\), obtaining
the newest parameter b.  Request the antecedent procedure in that
extension and use the original, now older, clause at b.  This yields
\(F(b)\).  The equality \(b=t\) has the persistent proof from the
atomic-specialization paragraph, even under the scopes later opened by
\(T_r\).  Apply \(T_r\) to obtain \(F(t)\), and close the naming
elimination.  This construction is uniform under embeddings.  It uses
\(T_r\) and the already established \(U_0\); the construction of
\(T_r\) used only ranks below r.  Hence the recursion is well founded.

\paragraph{Compiling arbitrary derivations.}
By structural recursion on an unrestricted proof, produce a persistent
procedure for its concluding sequent.  Refl and Step use the procedures
for their immediate subproofs.  At an unrestricted Use, its antecedent
procedure is already available by recursion on the proper subproof;
invoke \(U_r\) at the supplied numerical target.  At DE or DI, choose
fresh parameters in the requested larger scope, extend the context
embedding by the new clause, and invoke the previously constructed
procedure for the proper branch subproof.  Injective parameter renaming
preserves the formula root and the literal-zero witness test, so it
preserves the introduction mode.  Scope freshness and the clause-index
mapping are explicit in this procedure.

This compilation recurses on original proof height.  Its auxiliary T/U
procedures recurse on formula rank as described above.  There is no
recursion on an already generated proof of uncontrolled size.  Every
finite input proof therefore produces a finite restricted proof.
Evaluate the persistent output at the identity context embedding to
obtain the required derivation.
\end{proof}

\paragraph{Implementation significance.}
The matched complete checker replaces
\texttt{[u, index, target, premise]} by \texttt{[u, index, premise]}, removes
target-term validation from Use, and reverses the context-storage order.
The selected target is computed from the index and current scope.
For the incoming tagged uniform-Step implementation, the complete
normalized source changes from 2950 to 2938 bytes.  Significant tokens
remain 829, AST nodes change from 902 to 903, and informational gzip-9
size changes from 1005 to 999 bytes.  This is a mathematical removal of
the arbitrary specialization argument, with a small byte reduction;
it is not a strict improvement in every size coordinate or an absolute
minimality result.

\subsection{Ordinary realization and the interpretation composites}

Use the unchanged ordinary readings
\[
 \bigl(D_{x,y}B\bigr)^\circ
   =\exists x\bigl(B^\circ(x,0)\to\forall y\,B^\circ(x,y)\bigr),
 \qquad
 [A;y.B]^\circ=A^\circ\to\forall y\,B^\circ.
\]
Equality is actual numerical equality.  The underlying ordinary theory is
intuitionistic first-order logic with equality, the scalar iterator
equations, numerical coverage
\(\forall x(x=0\lor\exists y\,x=Sy)\), separation \(0=S0\to\bot\),
and induction for literal equations with side assumptions.  The inherited
uniform-Step equivalence identifies this ordinary theory with the
unrestricted five-rule clause presentation.  Its logical macro theorem
recovers all intuitionistic connectives and both quantifiers; its scalar
arithmetic bridge derives predecessor, reflection of equality, complete
quantifier-free induction, pairing, and a compiler for every primitive
recursive definition.

The new equivalence changes no term, formula, clause reading, or
arithmetical axiom.  It supplies proof translations in both directions on
exactly the same syntax: one direction includes Use$_*$ as the indicated
ordinary Use, and the other is the finite construction in the theorem.
Consequently it composes with the inherited whole-domain bi-interpretation
with full first-order iPRA.  The numerical domains in both directions are
the whole domain, equality remains actual equality, and the numerical
composite isomorphisms have graph \(x=y\).  The formula composites are
the inherited internally provable equivalences of the logical macros;
the present change adds an identity translation on formulas.  Thus
bi-interpretability and equiconsistency follow without a new arithmetic
bootstrap.  This paragraph relies on the inherited bridge theorem; the
new proof above establishes the specialization equivalence itself.

\paragraph{Scope of the novelty claim.}
The new item is the restriction of clause specialization to a target
computed solely from clause position and current scope, together with
the persistent-procedure argument that restores arbitrary specialization
and structural scope extension.  The iterator, five-rule architecture,
clausal contexts, scalar bridge, and arithmetic equality bootstrap are
inherited.  The reviewed source snapshot and the summaries of its filed
packages, together with the incoming reduced-clause, root-free-clause,
and uniform-Step packages, did not contain this deterministic target rule.
This is a scoped comparison, not a claim of priority over all possible
axiomatizations or a proof of global minimality.
